\section{Character-ratio exclusions for normal irreducible cores}
\label{sec:character-ratio}

We use the actual intermediate cover attached to an inducing block, as in
D-BI-01.  Thus an absolutely irreducible complex representation
$\tau:H\to\operatorname{GL}(V)$ of degree $e\ge2$ induces the original
representation, and its multiplicity local system has complex rank six.
Write $J=\tau(H)$, and discard precisely the branch points at which the
local $J$-monodromy is trivial.  No equality between the character field
of $\tau$ and the original imaginary quadratic field is assumed.

For a finite-order linear transformation $x$ put
\[
 \nu(x)=e-\max_{\lambda\in\mathbf C}\dim\ker(x-\lambda),
 \qquad d(x)=\operatorname{rank}(1-x).
\]
Thus $d(x)\ge\nu(x)$, whereas a nonidentity scalar has $\nu(x)=0$
and $d(x)=e$.  Keeping these two quantities distinct is essential below.

\begin{theorem}[D-BI-18: a character-ratio gate]
\label{D-BI-18}
Suppose $N\triangleleft J$ acts irreducibly on $V$, and for some
$0\le c<1$ its character $\chi$ satisfies
\[
          |\chi(x)|\le c e\qquad(x\in N\text{ nonscalar}).       \tag{BI.60}
\]
Every nonscalar $g\in J$ satisfies
\[
                         \nu(g)\ge\frac{e(1-c)}4.              \tag{BI.61}
\]
Consequently a period map of rank nine for the six-dimensional
multiplicity variation necessarily satisfies
\[
       12\left\lceil\frac{e(1-c)}4\right\rceil\le2e+6,
       \qquad e(1-3c)\le6.                                   \tag{BI.62}
\]
If $J\subset\mathbf C^*N$, the stronger necessary bounds are
\[
       12\left\lceil\frac{e(1-c)}2\right\rceil\le2e+6,
       \qquad e(4-6c)\le6.                                   \tag{BI.63}
\]
These conclusions allow inertia outside $N$, including inertia
inducing outer automorphisms of $N$.

In particular the following exclusions hold.
\begin{enumerate}
\item If $N$ is a finite quasisimple group of Lie type over
$\mathbf F_q$, its simple quotient is not alternating, and $q\ge47$,
then $e\in\{3,4,9\}$.  Thus every inducing degree $e\ge10$ is
excluded.  For the smaller field sizes in the following table, the
displayed upper bounds are necessary:
\[
\begin{array}{c|rrrrrr}
q&29&31&32&37&41&43\\ \hline
e\text{ at most}&87&45&33&21&15&15.
\end{array}                                                   \tag{BI.64}
\]
\item In the same Lie-type situation, if $J\subset\mathbf C^*N$
and $q\ge16$, then $e\le9$.  If $q\ge29$, no rank-nine source
exists at all.
\item If $N=ZE$, where $Z$ consists of scalars and $E$ is an
extraspecial group acting irreducibly, then $e\in\{3,4\}$.
In particular every such inducing degree $e\ge5$ is excluded.
\end{enumerate}
The same inequalities apply to an irreducibly acting central product
of commuting factors when each factor has a bound of the form
(BI.60): one may take the maximum of their character-ratio constants.
\end{theorem}

\begin{proof}
We first justify the moduli bound independently of any condition on
the field of definition.  Put $X=C/H$ and $Y=C/\ker\tau$.  The cover
$Y\to X$ has full deck group $J$.  Frobenius reciprocity identifies
the original multiplicity variation with the $\tau$-multiplicity
variation of $Y$, compatibly with its Hodge filtration.  For a fixed
topological type the possible covers $Y\to X$ are finite over the
moduli of the pointed curve $X$, so the actual period map factors
locally through that finite Hurwitz space.  Let $h=g(X)$, let $r$ be
the number of genuine branch points, and let their inertias be
$g_1,\ldots,g_r$.  Since $\tau$ is nontrivial and irreducible, the
Euler-characteristic formula is
\[
       6=2e(h-1)+D,\qquad D=\sum_{j=1}^r d(g_j),\qquad r\le D.
                                                                  \tag{BI.65}
\]
If $h=1$, then $D=6$ and the pointed-curve moduli have dimension at
most six.  If $h=2$, nonnegativity forces $e\le3$ and their dimension
is at most $3+D=9-2e\le5$.  If $h\ge3$ and $e\ge2$, then
$D=6-2e(h-1)<0$, which is impossible.  Therefore rank nine forces
$h=0$.  The dimension bound $r-3\ge9$ now gives
\[
                         r\ge12,\qquad D=2e+6,
                         \qquad e\ge3.                       \tag{BI.66}
\]
This is a necessary bound on a period map that already factors
through an actual intermediate Hurwitz space.  No dominance is
deduced from the dimension of that space.

Conjugate the finite group into a unitary group.  For $x\in N$
nonscalar, choose an eigenvalue $\lambda$ of largest multiplicity.
The trace of $\lambda^{-1}x$ is the sum of $e-\nu(x)$ ones and
$\nu(x)$ complex numbers of modulus one.  The reverse triangle
inequality, followed by (BI.60), gives
\[
                    e-2\nu(x)\le|\chi(x)|\le ce.
\]
Thus $\nu(x)\ge e(1-c)/2$.

Let $g\in J$ be nonscalar.  Since $N$ acts irreducibly, Schur's
lemma supplies $h\in N$ with $[g,h]\ne1$.  Normality gives
$x=[g,h]\in N$.  If $x$ is nonscalar, the preceding paragraph gives
$d(x)\ge\nu(x)\ge e(1-c)/2$.  If $x$ is a nonidentity scalar,
then $d(x)=e$, which is stronger.  For any nonzero scalar $\lambda$,
the commutator is unchanged when $g$ is replaced by
$A=\lambda^{-1}g$.  The elementary inequality
$\operatorname{rank}(1-AB)\le\operatorname{rank}(1-A)
+\operatorname{rank}(1-B)$ therefore gives
\[
 d([g,h])=d([A,h])
 \le d(A)+d(hA^{-1}h^{-1})=2d(A).
\]
Choose $\lambda$ of largest multiplicity in $g$ to get
$d([g,h])\le2\nu(g)$ and hence (BI.61).  This argument does not
require $N$ to be perfect.  A central but nonidentity commutator is
handled using its actual moving rank, not its largest-eigenspace
codimension.

Every genuine inertia has $d(g_j)\ge\lceil e(1-c)/4\rceil$:
this follows from (BI.61) for nonscalars, and from $d(g_j)=e$ for
nonidentity scalars.  Summing over the at least twelve entries in
(BI.66) proves the first inequality in (BI.62).  Dropping ceilings
and rearranging gives its second inequality.  If
$J\subset\mathbf C^*N$, a nonscalar $g=\alpha x$ has
$\nu(g)=\nu(x)\ge e(1-c)/2$ directly.  The identical summation
proves (BI.63).  These elementary support arguments are also the
characteristic-zero form of the arguments in
\cite[Definition~2.3 and Lemma~2.4, pp.~137--138;
Lemma~3.1, p.~139]{BI:DHP}; the proof above includes the scalar
adjustment needed for the invariant $\nu$.

For the Lie-type applications we use the published character bound
\cite[Theorem~4.1(2), p.~141]{BI:DHP}:
\[
 c(q)=\max\left\{\frac1{\sqrt q-1},\frac9q\right\}
                 \qquad(q>9).                              \tag{BI.67}
\]
That theorem attributes this simplified bound to Gluck
\cite{BI:Gluck1995}.  Its hypotheses are finite quasisimplicity,
Lie type over $\mathbf F_q$, and a nonalternating simple quotient.
It bounds every nontrivial irreducible complex character outside
its scalar locus.  Since $N\subset\operatorname{GL}(V)$ acts
faithfully and irreducibly, its scalar locus equals its center, so
it applies to exactly (BI.60).  We use the published formulation
in Duyan--Halasi--Podoski; no new proof of Gluck's character estimate
is claimed.

For $q\ge47$, both terms in (BI.67) are strictly less than $1/5$:
$47>45$ and $47>36$.  The continuous inequality in (BI.62) gives
$e<15$.  Its integer form excludes $5\le e\le8$, since its left
side is at least $24>2e+6$; it excludes $10\le e\le14$, since its
left side is at least $36>2e+6$.  Together with $e\ge3$, this
leaves precisely the necessary possibilities $3,4,9$.

For $q=29,31,32,37,41,43$ one has $c(q)=9/q$.  This comparison can
be checked without approximations by
$(q+9)^2\le81q$.  The continuous inequality gives, respectively,
$e\le87,46,38,22,17,16$.  Applying
$12\lceil e(q-9)/(4q)\rceil\le2e+6$ to these finite ranges
gives (BI.64).  Equivalently, set the ceiling equal to $k$ and
intersect the two rational intervals
$4q(k-1)/(q-9)<e\le4qk/(q-9)$ and $e\ge6k-3$.
The exact finite arithmetic is retained in
\texttt{check\_character\_ratio.py} and its JSON output.

For the scalar-only assertion, $q\ge16$ gives $c(q)\le9/16$,
and (BI.63) implies $e\le48/5$, hence $e\le9$.  If $q\ge29$,
then $9/q<1/3$ and $1/(\sqrt q-1)<1/3$, so (BI.63) gives
$e<3$, contrary to (BI.66).

For the extraspecial assertion, let $x\in E\setminus Z(E)$.
The nondegenerate commutator pairing of the extraspecial group
supplies $h\in E$ with $[h,x]=z\ne1$ in $Z(E)$.
Faithfulness and irreducibility make $z$ a nonidentity scalar.
As $hxh^{-1}=zx$, equality of traces gives
$\chi(x)=z\chi(x)$, hence $\chi(x)=0$.  Multiplication by a scalar
preserves vanishing, so every nonscalar element of $ZE$ has zero
trace.  Thus $c=0$.  Formula (BI.62) first gives $e\le6$, while
its integer form excludes $e=5,6$.  Since $e\ge3$, only $3,4$
remain.  This proves the stated exclusion in every degree at
least five, without a classification of all primitive groups.

Finally suppose a normal irreducible $N$ is generated by commuting
factors $N_1,\ldots,N_s$.  The irreducible representation is a
tensor product of irreducible factor representations $V_i$;
central identifications in a central product do not alter this
statement.  Its normalized character on $x=x_1\cdots x_s$ is
the product of the normalized characters of the $x_i$ on $V_i$.
A tensor product of invertible finite-order transformations is a
scalar only if every factor is a scalar, as is seen by comparing
their eigenvalue ratios.  For a nonscalar $x$, at least one
factor therefore has ratio at most $c_i$, and every other factor
has ratio at most one.  Taking $c=\max_i c_i$ proves (BI.60)
and completes the proof.
\end{proof}

\begin{corollary}[D-BI-18a: the essentially self-dual refinement]
\label{D-BI-18a}
In D-BI-18, suppose additionally that the original character field
is imaginary quadratic and that $J$ is essentially self-dual.
Then
\[
 2\lceil e/2\rceil+
 10\left\lceil\frac{e(1-c)}4\right\rceil\le2e+6,
 \qquad e(3-5c)\le12.                                      \tag{BI.68}
\]
If the original actual polarization is nonsplit, $\tau$ is primitive,
$e\ge9$, and $N$ is quasisimple of Lie type as above, then
$q\ge19$ is impossible.  These additional self-duality hypotheses
hold automatically when the irreducible representation of $N$
admits a nondegenerate invariant bilinear form.
\end{corollary}
\begin{proof}
Write $g^{\mathsf T}Bg=\alpha(g)B$ for the similitude character of
$J$.  If $\alpha$ were trivial, the character of $J$ would be real.
Induction of characters includes the original imaginary quadratic
field in the character field of $\tau$, so this is impossible.
The genuine inertias generate $J$ and have product one, hence at
least two have $\alpha(g_j)\ne1$.  Their fixed spaces are totally
isotropic for $B$, so each moves at least $\lceil e/2\rceil$
dimensions.  At least ten other inertias remain, and D-BI-18
applies to each.  This proves the integer inequality in (BI.68);
discarding ceilings gives its continuous form.

For $q\ge19$, (BI.67) gives $c(q)<1/2$, so (BI.68) implies
$e<24$.  D-BI-09 excludes primitive essentially self-dual inducing
degrees $9\le e\le23$ for a dominant nonsplit sixfold source,
including larger inducing character fields.  Together these two
statements prove the claimed exclusion.

Finally let $B$ be an $N$-invariant nondegenerate bilinear form.
For every $g\in J$, the form $g^{\mathsf T}Bg$ is again
$N$-invariant, because $g$ normalizes $N$.  Irreducibility makes
the space of such forms one-dimensional by Schur's lemma.
Therefore $g^{\mathsf T}Bg=\alpha(g)B$ for a character $\alpha$,
which proves the last assertion.
\end{proof}

\paragraph{Scope.}
The result applies to arbitrary inducing character fields and does
not require the whole inducing group to be quasisimple.  It does
require the stated normal subgroup to act irreducibly.  The
existence of such a subgroup has not been assumed for every
primitive group.  A reducible normal core, factors without an
adequate character-ratio bound, and small defining fields remain
separate cases.  The alternating groups are deliberately excluded
from (BI.67); their character ratios have no uniform upper bound
strictly below one, as also noted in
\cite[Remark~4.2, p.~141]{BI:DHP}.  None of these arguments is a
claim of a universal quarter-dimensional moving-rank theorem for
all finite primitive complex linear groups.
